%=============================================================================
% 学位论文选题依据与意义 + 创新点
%=============================================================================

\section*{一、选题依据与意义}

中微子作为连接粒子物理与宇宙学的关键粒子，其非零质量与振荡现象是迄今唯一具有坚实实验证据的超越标准模型的新物理信号。自 1998 年 Super-Kamiokande 实验首次发现大气中微子振荡以来，经过 SNO、KamLAND、Daya Bay 等一系列里程碑式实验的推进，三代中微子混合框架已基本确立，中微子物理正从"现象发现"迈向"精密测量"时代。在这一新阶段，实现振荡参数的亚百分比量级精确测量具有深远的物理意义：一方面，对太阳振荡参数 $\sin^2\theta_{12}$ 和 $\Delta m^2_{21}$ 的高精度测定，是未来无中微子双贝塔衰变（$0\nu\beta\beta$）实验确定中微子有效质量上限及质量序的前提条件，直接关系到判定中微子究竟是马约拉纳（Majorana）粒子还是狄拉克（Dirac）粒子这一根本问题；另一方面，大气振荡参数 $\Delta m^2_{31}$ 的精确测量及质量序的判定，对完善大统一理论和理解宇宙物质--反物质不对称性至关重要。

江门中微子实验（JUNO）作为国际领先的中基线反应堆中微子实验，依托地下 700 米处 20 千吨液体闪烁体探测器，在 52.5 公里基线处探测反应堆反中微子，旨在判定中微子质量序并实现振荡参数的亚百分比量级测量。JUNO 早期仅 59.1 天的运行数据便将太阳振荡参数的测量精度较全球既有结果提升了 1.6 倍，有力验证了探测器卓越的能量分辨率与大规模工程设计的可行性。然而，实验早期运行面临三大关键挑战：其一，低统计量条件下能谱分箱过细会导致高斯近似失效，引入可观的参数估计偏差；其二，反应堆能谱模型的不确定性在精细分段下可能引入难以忽略的系统误差；其三，在参数空间存在多重局部极小值的情况下，Wilks 定理的渐近近似可能不再成立，直接影响置信区间的频率学覆盖率。

针对上述挑战，迫切需要开发高性能拟合框架以支撑海量伪实验分析，系统研究能谱分箱与分段策略以平衡灵敏度与系统误差控制，并运用严格的频率学统计方法确保物理结论的稳健性。本文正是围绕反应堆中微子振荡分析的完整物理链条，从数值算法优化、能谱分段策略与严格统计方法三个维度开展了深入研究，为 JUNO 实验早期物理结果的产出提供了关键的方法学支撑与定量依据。


\section*{二、创新点}

\subsection*{创新点 1：高性能反应堆中微子张量化拟合框架}

针对 Feldman--Cousins 覆盖率检验所需的百万量级拟合对计算资源的极高需求，本文设计并实现了一套基于张量化思想的前向折叠反应堆能谱拟合框架。该框架将反应堆中微子能谱生成、逆贝塔衰变（IBD）微分截面计算、沉积能谱映射及探测器响应建模统一构建为张量运算链，并引入三项核心优化技术：结果缓存机制消除重复中间计算、带状稀疏矩阵表示利用响应矩阵的局域性降低内存与计算开销、分块张量化构造与分段卷积优化最小化内存搬运。基准测试表明，该框架在保持数值精度（相对误差低于 $10^{-6}$）的前提下，较传统实现获得约一个数量级的加速。该拟合器已被 JUNO 合作组采纳为基础振荡分析工具，为大规模参数扫描与频率学覆盖率研究提供了可行的计算平台。相关成果已发表于 \textit{Eur.\ Phys.\ J.\ C}（2025）。

\subsection*{创新点 2：低统计量下最优能谱分箱与分段策略}

针对 JUNO 实验早期低统计量条件下拟合偏差与反应堆模型依赖性问题，本文通过大规模伪实验系统评估了重建能谱的分箱（binning）策略与反应堆能谱的分段（segmentation）方案。研究发现，当重建能谱箱宽达到 100~keV 及以上时，由非高斯统计涨落引入的振荡参数偏差趋于稳定且处于可接受范围。在 JUNO--大亚湾联合分析框架下，本文进一步提出将反应堆能谱分段宽度设定为 300~keV，该策略在抑制反应堆模型依赖（系统误差贡献低于 0.3\%）与保留物理灵敏度之间取得了有效平衡，避免了因分段过细而超出近点实验分辨能力所产生的非物理结构。上述分箱与分段策略已被 JUNO 合作组采纳，并应用于首篇物理结果论文的振荡分析中。

\subsection*{创新点 3：基于 Feldman--Cousins 方法的统计稳健性推断}

针对 Wilks 定理在低统计量或似然函数结构复杂情形下的适用性问题，本文运用 Feldman--Cousins（FC）方法对中微子振荡参数开展了大规模频率学覆盖率检验。结果表明，太阳振荡参数 $\Delta m^2_{21}$ 与 $\sin^2\theta_{12}$ 的二维 $\Delta\chi^2$ 分布与自由度为 2 的卡方分布一致，Wilks 近似成立；然而对大气振荡参数 $\Delta m^2_{31}$，经验 $1\sigma$ 临界值约为 1.57，显著高于 Wilks 定理预期的标准值 1.0，揭示了在当前低统计量下直接使用渐近阈值将导致覆盖率不足。基于 FC 方法，本文给出了 JUNO 早期运行 50 天统计量下 $\Delta m^2_{31}$ 约 1.4\% 的测量精度。该结果为合作组在首篇物理结果论文中的统计处理方案提供了严格的频率学依据。
